\begin{table*}[t!]
\centering
%\small
%\rowcolors{2}{Para}{white}
\arrayrulecolor{black}
\resizebox{\textwidth}{!}{
\begin{tabular}{
>{\raggedright\arraybackslash}m{2cm}
>{\raggedright\arraybackslash}m{2cm}
>{\raggedright\arraybackslash}m{1cm}
>{\raggedright\arraybackslash}m{1cm}
>{\raggedright\arraybackslash}m{1cm}
>{\raggedright\arraybackslash}m{1.5cm}
>{\raggedright\arraybackslash}m{2cm}
>{\raggedright\arraybackslash}m{5cm}
>{\raggedright\arraybackslash}m{1cm}
}
\toprule
 \textbf{Method} 
& \textbf{Type} 
& \textbf{TF} 
& \textbf{TB}
& \textbf{TS} 
& \textbf{Domain} 
& \textbf{LMs} 
& \textbf{Main Advancement}
& \textbf{Year} \\
\toprule

\textbf{HippoRAG 2} \newline\cite{gutierrez2025rag}
& 
& \usym{2717}
& \usym{2717}
& Task-Free
& Question Answering
& 
& Employs a training objective that \textbf{minimizes the Kullback-Leibler (KL) divergence} between the predictions of the original model and target model. 
& 2025 \\ 

\textbf{SELF-PARAM} \newline\cite{wangself}
& Regularization-based Learning
& \usym{2713}
& \usym{2713}
& Task-Free
& Question Answering
& Llama-3.3-70B-Instruct
& Enhances Personalized PageRank-based retrieval with deeper passage integration and online LLM usage, achieving superior performance on factual, associative, and sense-making memory tasks. 
& 2025 \\ 

\textbf{MBPA++} \newline\cite{wang2024towards}
& Replay-based
& \usym{2717}
& \usym{2717}
& CIL
& None
& REPLAY, MBPA
& Maintains a small, randomly selected subset (as low as 1\%) of past examples in memory to achieve performance comparable to larger memory sizes.
& 2025 \\ 

\textbf{LSCS} \newline\cite{wang2024towards}
& Interactive Learning
& \usym{2717}
& \usym{2717}
& CIL
& Abstracting/ Merging/ Retrieval
& /
& Integrates \textbf{multiple storage mechanisms} to achieve abstraction, experience merging, and long-term retention with accurate recall.
& 2025 \\ 

\textbf{TaSL} \newline\cite{feng2024tasl}
& Regularization-based Learning
& \usym{2717}
& \usym{2717}
& TIL
& Dialogue System
& T5, Llama-7B
& Parameter-level task skill localization and consolidation enabling knowledge transfer \textbf{without memory replay}.
& 2024 \\ 

\textbf{EMP} \newline\cite{liu2022incremental}
& Replay-based
& \usym{2717}
& \usym{2717}
& CLI
& Event Detection
& BERT-ED, KCN
& Designs \textbf{continuous prompts} associated with each event type.
& 2023 \\ 

\textbf{UDIL} \newline\cite{shi2023unified}
& Interactive Learning
& \usym{2717}
& \usym{2713}
& DLI
& Event Detection
& oEWC, SI, LwF, A-GEM, CLS-ER, ESM, etc.
& Introduces \textbf{adaptive coefficients} optimized during training to achieve tighter generalization error bounds and improved performance across domains.
& 2023 \\ 

\textbf{DSI++} \newline\cite{mehta2022dsi++}
& Replay-based
& \usym{2717}
& \usym{2713}
& TIL
& Information Retrieval
& T5
& Enables \textbf{continual document indexing} while retaining query performance on old and new data.
& 2022 \\ 

\textbf{MRDC} \newline\cite{wang2022memory}
& Replay-based
& \usym{2717}
& \usym{2713}
& CIL
& Object Detection
& LUCIR, PODNet
& Enhances \textbf{memory replay by compressing data}, balancing sample quality and quantity for continual learning.
& 2022 \\ 

\bottomrule
\end{tabular}
}
\caption{Overview of methods for \textbf{parametric memory modification in continual learning}. "TB" denotes whether task boundaries exist. "TS" denotes task settings including TIL (Task Incremental Learning), CIL (Class Incremental Learning), DIL (Domain Incremental Learning), and Task-Free.}
\label{tab:methods-continual_learning}
\end{table*}
